\subsection{LLMs Cannot Reliably Identify and Reason About Security Vulnerabilities (Yet?): A Comprehensive Evaluation, Framework, and Benchmarks}
\label{sec:vulnerability-detection}

\subsubsection{Motivation}
Ullah et al.~\cite{ullah2024secllmholmes} ask whether chat-based LLMs can detect software vulnerabilities and explain their decisions reliably. An alert is hard to act on when its explanation directs the reviewer to the wrong cause.

\subsubsection{Methodology}
SecLLMHolmes evaluated eight models across 228 code scenarios, including 48 handcrafted examples, 30 cases drawn from 15 real-world CVEs, and 150 modified examples. The benchmark asked about a specified weakness and included vulnerable and patched versions. The framework scored verdicts and explanations separately and tested how results changed across repeated runs, prompts, and code edits~\cite[Secs.~3--4]{ullah2024secllmholmes}.

\subsubsection{Findings}
GPT-4 correctly classified 43 of 48 handcrafted cases under several prompts. Some correct verdicts came with incorrect explanations. In tests using edits that preserved the target vulnerability status, the authors highlighted error rates of 26\% for PaLM~2 and 17\% for GPT-4. On real-world cases, the models missed vulnerabilities and raised false alarms on patched code~\cite[Secs.~4.3--4.8]{ullah2024secllmholmes}.

\subsubsection{Conclusions}
The small real-world sample, the specified target weakness, and the tested model versions limit how broadly the results apply. A reviewer can use a model alert to decide what to inspect and then verify the claim against the code. The study evaluates detection in existing code, a task that also matters when reviewing code generated by an assistant~\cite{ullah2024secllmholmes}.
